% SLIM-ARC: 学术报告 - 核心设计
\section{核心设计}

\subsection{整体架构}

SLIM-ARC 采用分层架构设计，如图\ref{fig:architecture}所示，自底向上分为三层。这一架构的核心思路是：\textbf{利用操作系统虚拟内存机制作为 MoE 稀疏性的天然接口}，避免在用户态重新实现复杂的内存管理（如 FlexInfer~\cite{flexinfer2025}的 Direct I/O + 手动 buffer 管理），而是通过 \texttt{madvise} 等标准 POSIX 接口~\cite{posix_madvise}控制内核的页面调度行为。

\begin{enumerate}
    \item \textbf{内核协同层}（mmap + madvise）：利用操作系统虚拟内存机制，将 GGUF 模型文件通过 mmap 映射到进程虚拟地址空间（45\,GB VSZ），通过 \texttt{posix\_madvise} 控制内核的页面预取行为，实现按需分页。此层的设计灵感来自 llama.cpp~\cite{llamacpp2024}的 mmap 机制，但针对 MoE 稀疏访存模式进行了策略调整。
    \item \textbf{运行时调度层}（prefetch\_scheduler + unified\_io\_scheduler + MoE Router Hook）：基于运行时阶段（Prefill/Decode）和 MoE Router 输出，动态调度权重预取、专家选择性加载和 KV Cache 管理。此层借鉴了 FlexInfer 的层感知预取和 MobileMoE~\cite{mobilemoe2024}的 Router Predictor 思路。
    \item \textbf{量化优化层}（IQ4\_XS + KV q4\_0 + FlashAttention）：通过模型权重的低比特量化（IQ4\_XS，45$\to$40\,GB）、KV Cache 量化（f16$\to$q4\_0）和 FlashAttention~\cite{flashattention2022}的 IO-aware tiling 融合，从算法层面降低物理内存占用和 I/O 带宽需求。
\end{enumerate}

% Figure: SLIM-ARC 整体架构图（方框占位）
\begin{figure}[H]
    \centering
    \fbox{\parbox{0.9\textwidth}{\centering
    \vspace{4em}
    \textbf{[图占位：SLIM-ARC 三层架构示意图]}\\[1em]
    \textit{顶层：量化优化层} —— IQ4\_XS 权重 (40\,GB) $\to$ KV q4\_0 $\to$ FlashAttention (IO-aware tiling)\\[0.5em]
    \textit{中层：运行时调度层} —— prefetch\_scheduler + MoE Router Hook + unified\_io\_scheduler\\[0.5em]
    \textit{底层：内核协同层} —— mmap (45\,GB VSZ) + posix\_madvise (MADV\_RANDOM / WILLNEED / DONTNEED)\\[0.5em]
    $\uparrow\downarrow$ 数据流：SSD (GGUF) $\to$ Page Cache $\to$ CPU Registers\\
    \vspace{4em}
    }}
    \caption{SLIM-ARC 整体架构。三层自底向上：内核协同层通过 mmap 将 45\,GB GGUF 映射到虚拟地址空间，用 madvise 控制页面预取策略；运行时调度层基于 MoE Router 输出和计算阶段动态调度预取；量化优化层从算法侧降低内存和带宽需求。MADV\_RANDOM 是核心：它让内核的 demand paging 与 MoE 2\% 的稀疏激活率匹配，仅加载被 Router 选中的专家权重。}
    \label{fig:architecture}
\end{figure}

\subsection{mmap + MADV\_RANDOM 按需加载机制}

\subsubsection{设计动机}

在传统的 mmap 加载方式下，llama.cpp~\cite{llamacpp2024}使用 \texttt{madvise(WILLNEED)} 提示内核进行顺序预读。对于小模型（如 4B），这能加速 prefill 阶段的顺序访问。但对于 80B 大模型，内核的 readahead 策略会预读大量未使用的页面，导致 8\,GB 环境下的物理内存耗尽和 OOM。

关键观察是：MoE 模型在 decode 阶段的访存具有高度稀疏性。Qwen3-Next-80B~\cite{qwen3next2024}每个 token 仅激活 10/512 个专家，其余 502 个专家的权重完全不需要加载。如果能阻止内核预读这些未激活的页面，物理内存占用可从 45\,GB 降至接近实际访问量（约 2\,GB）。这一思路与 PowerInfer~\cite{powerinfer2024}利用 ReLU 激活稀疏性的原理类似，但 SLIM-ARC 利用的是 MoE 路由稀疏性，且无需修改模型结构。

\subsubsection{机制设计}

在模型加载阶段（\texttt{init\_mappings}），对大于 6\,GB 的 mmap 区域调用 \texttt{posix\_madvise(MADV\_RANDOM)}~\cite{posix_madvise}：

\begin{equation}
    \texttt{posix\_madvise(addr, size, POSIX\_MADV\_RANDOM)}
\end{equation}

\texttt{MADV\_RANDOM} 的语义是告诉内核该区域的访问模式是随机的，内核应禁用 readahead，仅在页面被实际访问时触发缺页（page fault）加载。

对于 decode 阶段，每个 token 的前向传播只访问当前层的权重（非 MoE 部分）和被激活的专家权重。MADV\_RANDOM 使得这些访问以 4\,KB 页粒度按需加载，未访问的专家权重完全不占用物理内存。

\begin{figure}[H]
    \centering
    \fbox{\parbox{0.9\textwidth}{\centering
    \vspace{3em}
    \textbf{[图占位：MADV\_RANDOM 按需加载机制对比]}\\[1em]
    \textit{左：默认 WILLNEED} —— 内核顺序预读所有 512 专家页面 $\to$ 45\,GB RSS $\to$ OOM\\[0.5em]
    \textit{右：MADV\_RANDOM} —— 仅 page fault 加载 10 个激活专家 $\to$ $\sim$2\,GB RSS\\
    \vspace{3em}
    }}
    \caption{MADV\_RANDOM 按需加载机制。默认 WILLNEED 策略下，内核预读整个 mmap 区域（45\,GB），导致 8\,GB 环境 OOM；MADV\_RANDOM 禁用预读，仅在被 Router 激活的专家页面被访问时触发 page fault，物理内存占用降至 $\sim$2\,GB（10/512 $\times$ 45\,GB $\approx$ 0.88\,GB + 非专家权重）。}
    \label{fig:madv_random}
\end{figure}

\subsubsection{prefill/decode 的 tradeoff}

MADV\_RANDOM 在 decode 阶段带来显著加速，但对 prefill 阶段有害。prefill 需要顺序读取所有层的全部权重，MADV\_RANDOM 阻止了内核的顺序预读，导致每个页面都要同步缺页。实测中 prefill 速度下降约 57\%，但 decode 速度提升 200--850\%。

由于交互式推理的主要瓶颈在 decode（逐 token 生成），这一 tradeoff 对实际应用是有利的。系统通过环境变量 \texttt{SLIM\_ARC\_NO\_MADV\_RANDOM=1} 提供了关闭选项，供 prefill-heavy 场景使用。

\subsection{禁用 GGML\_CPU\_REPACK}

\subsubsection{问题发现}

在集成 MADV\_RANDOM 后，80B 模型在 8\,GB 环境下仍然 OOM。通过 \texttt{dmesg} 分析发现，进程的匿名物理内存（\texttt{anon-rss}）达到 8.3\,GB，远超模型权重本身。根因在于 llama.cpp 的 CPU 后端默认启用了 \texttt{GGML\_CPU\_REPACK}（CMake 选项 \texttt{-DGGML\_CPU\_REPACK=ON}），该机制将 Q4\_K 权重重打包为 \texttt{q4\_K\_8x8} 格式以优化 GEMM 的 SIMD 访问模式，但这会在内存中生成一份完整的重打包副本，导致匿名内存翻倍（45\,GB $\to$ 90\,GB）。

\subsubsection{解决方案}

通过 CMake 编译选项 \texttt{-DGGML\_CPU\_REPACK=OFF} 禁用重打包，CPU 后端直接使用 mmap 区域中的原始 Q4\_K 权重进行计算。实测验证，禁用 repack 后：80B 模型在 8\,GB 环境下不再 OOM（\texttt{anon-rss} 降至约 1\,GB）；热缓存下 decode 性能无损失（repack 的 SIMD 优化对 memory-bound 的 decode 无显著帮助）；prefill 性能略有下降，但被 MADV\_RANDOM 的 prefill 代价所掩盖。

\subsection{层感知预取调度器}

\subsubsection{架构设计}

\texttt{prefetch\_scheduler} 在模型加载时注册所有权重张量的 mmap 地址、大小和层级信息。在每次 \texttt{graph\_compute} 调用时，遍历计算图节点提取引用的层级，对当前层 + 预取窗口内的后续层调用 \texttt{madvise(WILLNEED)} 提示内核异步预读。Prefill 阶段使用较大的窗口（默认 4 层），因为 prefill 是计算密集型，I/O 可被计算掩盖；Decode 阶段使用较小窗口（1 层），因为 decode 是访存密集型，精确预取更重要。其调度循环如算法\ref{alg:prefetch}所示。

\begin{lstlisting}[caption={prefetch\_scheduler 调度循环伪代码}, label={alg:prefetch}, language=Python, basicstyle=\footnotesize\ttfamily]
# Called at each graph_compute(graph, batched)
def prefetch_tick(graph, batched):
    min_layer, max_layer = extract_layer_range(graph)
    phase = PREFILL if batched else DECODE
    window = window_prefill if phase == PREFILL else window_decode
    
    # 1. Notify current layer (may trigger async WILLNEED for window layers)
    for l in range(min_layer, min(max_layer + 1, min_layer + window + 1)):
        if l > max_layer:
            # Speculative prefetch of future layers
            for tensor in tensors_by_layer[l]:
                posix_madvise(tensor.addr, tensor.size, POSIX_MADV_WILLNEED)
    
    # 2. MoE expert prefetch (if router cache available)
    for l in range(min_layer, max_layer + 1):
        experts = router_expert_cache.get(l - 1)  # predict from prev layer
        if experts:
            prefetch_experts(l, experts)
    
    # 3. Evict old layers (DONTNEED) to free page cache
    for l in range(0, min_layer - window):
        for tensor in tensors_by_layer[l]:
            if tensor.size > 6GB:  # only large tensors
                posix_madvise(tensor.addr, tensor.size, POSIX_MADV_DONTNEED)
\end{lstlisting}

\subsubsection{MoE Router Hook 与跨层专家预测}

利用 llama.cpp 计算图中 \texttt{ffn\_moe\_topk} 张量包含 Router 输出的 top-k expert IDs 这一特性，SLIM-ARC 实现了跨层专家预测：在每次 \texttt{graph\_compute} 完成后，遍历计算图节点查找名为 \texttt{ffn\_moe\_topk} 的张量，读取其 I32 数据（形状为 \texttt{[n\_expert\_used, n\_tokens]}）提取每个 token 的激活专家 ID，缓存到 \texttt{cache\_router\_experts(layer, expert\_ids)}，在下一次 \texttt{graph\_compute} 开始时用上一层的 router 输出预测当前层的激活专家，对预测的专家在合并 3D 张量中的对应子区域调用 \texttt{madvise(WILLNEED)}。

\begin{figure}[H]
    \centering
    \fbox{\parbox{0.9\textwidth}{\centering
    \vspace{3em}
    \textbf{[图占位：MoE Router Hook 跨层专家预测流程]}\\[1em]
    Layer N 计算完成 $\to$ 提取 \texttt{ffn\_moe\_topk} $\to$ 缓存 expert IDs\\
    $\downarrow$\\
    Layer N+1 开始 $\to$ 用 Layer N 的 experts 预测 $\to$ \texttt{madvise(WILLNEED)} 子区域\\
    $\downarrow$\\
    内核异步预读 10 个专家权重（而非全部 512 个）\\
    \vspace{3em}
    }}
    \caption{MoE Router Hook 跨层专家预测。llama.cpp 计算图中的 \texttt{ffn\_moe\_topk} 张量包含 Router 输出的 top-k expert IDs。SLIM-ARC 在每层计算后提取该张量，缓存 expert IDs，用于预测下一层的激活专家并提前 \texttt{madvise(WILLNEED)}。由于 Qwen3-Next-80B 的专家权重以 \texttt{[dim, dim, 512]} 的 3D 合并张量存储，每个专家占 \texttt{total\_size / 512} 的连续子区域，通过计算偏移量 \texttt{expert\_addr = base + expert\_id $\times$ per\_expert\_size} 可精确预取。}
    \label{fig:moe_router_hook}
\end{figure}

\subsection{KV Cache 量化}

将 KV Cache 的数据类型从默认的 F16 改为 Q4\_0，使 KV Cache 的内存占用减半。对于 80B 模型在 16K 上下文下，KV Cache 从 6\,GB 降至 3\,GB，为权重缓存腾出更多物理内存。实测 decode 性能提升约 14\%。这一设计参考了 KVQuant~\cite{kvquant2024}和 KIVI~\cite{kivi2024}的 4-bit KV 量化方案，但使用 llama.cpp 原生的 Q4\_0 格式以避免额外计算开销。综述~\cite{survey2026}Sec 3.3 指出，4-bit KV 量化在多数任务上的精度损失可控，SLIM-ARC 的 GSM8K 实验验证了这一点。

\subsection{统一 I/O 带宽预算调度器}

\texttt{unified\_io\_scheduler} 根据运行时阶段（PREFILL\_SHORT, PREFILL\_LONG, DECODE\_SHORT, DECODE\_LONG, MOE\_DECODE）动态分配 I/O 带宽预算，协调权重预取、KV Cache 换页和专家预取三路 I/O 竞争。其设计灵感来自 DUAL-BLADE~\cite{dualblade2024}的"带宽预算"思路和综述~\cite{survey2026}Sec 4.3 的"elastic memory hierarchy"建议。

\begin{table}[H]
    \centering
    \caption{统一 I/O 带宽分配比例}
    \label{tab:budget_ratios}
    \begin{tabular}{lccc}
        \toprule
        \textbf{阶段} & \textbf{权重} & \textbf{KV Cache} & \textbf{专家} \\
        \midrule
        PREFILL\_SHORT & 60\% & 10\% & 30\% \\
        PREFILL\_LONG  & 50\% & 20\% & 30\% \\
        DECODE\_SHORT  & 70\% & 20\% & 10\% \\
        DECODE\_LONG   & 30\% & 60\% & 10\% \\
        \bottomrule
    \end{tabular}
\end{table}

\subsection{FlashAttention 集成}

受综述~\cite{survey2026}Sec 4.1 "FlashAttention~\cite{flashattention2022} fuses attention IO-aware" 启发，SLIM-ARC 启用了 llama.cpp 内置的 FlashAttention 实现（\texttt{-fa auto} 参数）。FlashAttention 通过 IO-aware tiling 融合，将 attention 的 QK$^T$ 和 softmax(P)V 合并为单个 kernel，避免中间矩阵的显存读写，对 memory-bound 的 decode 阶段效果尤为显著。实测 80B decode 吞吐量提升 71.4\%（3.01 $\to$ 5.16 t/s），零代码改动、零精度损失。

\subsection{StreamingLLM KV Eviction}

受综述~\cite{survey2026}Sec 4.3 "combine KV compression with eviction policies (e.g., window + sink tokens)" 启发，SLIM-ARC 实现了 StreamingLLM~\cite{streamingllm2024}式的 KV Cache 驱逐机制。在每次 decode step 后检查 KV 序列长度，若超过阈值 \texttt{sink + window}，则驱逐区间 \texttt{[sink, seq\_len - window)} 内的 token，仅保留 attention sinks（前 $N=4$ 个 token，稳定 softmax 分母）和 sliding window（最近 $W$ 个 token）。其算法如算法\ref{alg:eviction}所示。

\begin{lstlisting}[caption={StreamingLLM KV Eviction 伪代码}, label={alg:eviction}, language=Python, basicstyle=\footnotesize\ttfamily]
# Called after each decode step (batched=False)
def maybe_evict(seq_len, rm_fn):
    if not enabled: return
    threshold = sink_size + window_size  # e.g., 4 + 1024 = 1028
    if seq_len <= threshold: return
    
    # Evict tokens in [sink, seq_len - window), keeping:
    #   [0, sink) as attention sinks
    #   [seq_len - window, seq_len) as recent sliding window
    p0 = max(last_evicted_pos, sink_size)
    p1 = seq_len - window_size
    if p1 > p0:
        rm_fn(seq_id=0, p0=p0, p1=p1)  # llama_memory_seq_rm
        last_evicted_pos = p1
        # KV memory: O(context_length) -> O(sink + window)
\end{lstlisting}

\begin{figure}[H]
    \centering
    \fbox{\parbox{0.9\textwidth}{\centering
    \vspace{3em}
    \textbf{[图占位：StreamingLLM KV Eviction 示意图]}\\[1em]
    KV Cache: \texttt{[sink(4)][----evicted----][window(1024)]}\\
    $\uparrow$ 永久保留 $\quad\quad\quad$ $\uparrow$ 驱逐 $\quad\quad\quad$ $\uparrow$ 滑动保留\\
    内存: $O(L)$ $\to$ $O(\text{sink} + \text{window})$ = $O(1)$\\
    \vspace{3em}
    }}
    \caption{StreamingLLM KV Eviction 机制。保留前 4 个 token 作为 attention sink（稳定 softmax 分母，避免 PPL 暴增）和最近 $W$ 个 token 作为 sliding window，驱逐中间所有 token。KV Cache 内存从 $O(\text{context\_length})$ 降至 $O(\text{sink} + \text{window})$，实现恒定内存下的长上下文推理。该机制通过 \texttt{llama\_memory\_seq\_rm} 接口实现，无需修改 attention kernel。}
    \label{fig:kv_eviction}
\end{figure}
